\section{有理逼近与 Stern-Brocot 树}


这一部分来自课堂后半段讨论，手写三页中没有完整记录。整理时补齐了分母引理、终止性和唯一性的证明。

\subsection{中项搜索}


给定 $x>0$，从端点 $0/1$ 与 $1/0$ 出发，其中 $1/0$ 只是表示 $+\infty$ 的哨兵，不作通常的除法。对端点 $a/b<c/d$，取中项

\[
\frac{a+c}{b+d}.
\]

若它等于 $x$ 就停止；若大于 $x$，用它替换右端点；若小于 $x$，用它替换左端点。中项是分子、分母分别相加，不是两个分数的算术平均。

例如 $x=(\sqrt5-1)/2\approx0.618$ 的搜索路径依次比较 $1/1,1/2,2/3,3/5,5/8,8/13,\ldots$。这是黄金比例的倒数；通常所称黄金比例 $(1+\sqrt5)/2\approx1.618$。

\subsection{行列式不变量与最简性}


始终有

\[
bc-ad=1.
\]

初始时成立；替换任一端点后仍成立，例如 $b(a+c)-a(b+d)=bc-ad$。用矩阵写成

\[
M=\begin{pmatrix}a&c\\b&d\end{pmatrix},\qquad\det M=-1.
\]

保留左区间对应右乘 $L$，保留右区间对应右乘 $R$，其中

\[
L=\begin{pmatrix}1&1\\0&1\end{pmatrix},\qquad
R=\begin{pmatrix}1&0\\1&1\end{pmatrix}.
\]

二者行列式都为 $1$，所以不变量保持。中项也是最简分数：若 $g\mid a+c$ 且 $g\mid b+d$，则 $g$ 整除 $b(a+c)-a(b+d)=1$，故 $g=1$。

\subsection{相邻端点间的分母下界}


设 $b,d>0$、$a/b<c/d$ 且 $bc-ad=1$。若整数 $p,q$ 满足 $q>0$ 且

\[
\frac ab<\frac pq<\frac cd,
\]

则 $q\ge b+d$。证明：交叉相乘得到正整数

\[
pb-aq\ge1,\qquad cq-pd\ge1.
\]

分别乘 $d,b$ 再相加：

\[
q(bc-ad)=d(pb-aq)+b(cq-pd)\ge b+d.
\]

由于 $b,d>0$，等号要求两个正整数都为 $1$，解得 $p=a+c,q=b+d$。所以中项是开区间内唯一具有最小分母 $b+d$ 的分数。

注意：这个证明不能在 $d=0$ 时直接使用；下一步先处理最初的无穷端点。

\subsection{每个正有理数都出现且只出现一次}


固定最简正有理数 $x=p/q$。若最初一直保留右区间，中项会依次为 $1,2,3,\ldots$。因为 $x$ 有限，有限步内就会命中 $x$ 或得到两个分母都为正的端点。

假设此后仍永不命中，那么 $x$ 始终严格夹在端点之间。每一步分母从 $(b,d)$ 变为 $(b,b+d)$ 或 $(b+d,d)$，分母和至少增加 $1$，最终必有 $b+d>q$。但分母引理又要求 $q\ge b+d$，矛盾。因此搜索必定命中。

唯一性来自区间分割：每次新中项严格位于当前区间内部，其左右两个子树分别落在两个不相交的开区间中；已出现的中项作为边界，不会再在子树内部生成。故每个正有理数在整棵树中恰出现一次。固定一个目标的搜索路径只是一条路径，不会遍历所有有理数。

反向合并连续的左步或右步对应整数倍的列减法，与辗转相除和连分数有关；这里不据此把每个中项都称作连分数的收敛分数。

\subsection{“最佳逼近”必须说明比较范围}


分母引理保证的是\textbf{两个相邻端点之间}中项的分母最小，不保证搜索路径上的每个点都对“绝对误差、分母”同时最优。

反例：$x=51/100$ 的路径包含 $1/2$ 和 $2/3$，而

\[
\left|\frac{51}{100}-\frac12\right|=\frac1{100}
<\frac{47}{300}=\left|\frac{51}{100}-\frac23\right|,
\]

且 $2<3$。所以 $1/2$ 在这两个指标上都优于 $2/3$，原稿“路径上的每个分数都在帕累托前沿”应删除。若讨论单侧逼近或其他误差指标，需要另行陈述定理。

对给定的 Farey 序列 $F_N$，已在序列中的最简分数 $a/b<c/d$ 相邻，还需要 $bc-ad=1$ 和 $b+d>N$；只写行列式为 $1$ 没有包含阶数 $N$ 的限制。
